% Options for packages loaded elsewhere
\PassOptionsToPackage{unicode}{hyperref}
\PassOptionsToPackage{hyphens}{url}
%
\documentclass[
  11pt,
]{article}
\usepackage{lmodern}
\usepackage{amssymb,amsmath}
\usepackage{ifxetex,ifluatex}
\ifnum 0\ifxetex 1\fi\ifluatex 1\fi=0 % if pdftex
  \usepackage[T1]{fontenc}
  \usepackage[utf8]{inputenc}
  \usepackage{textcomp} % provide euro and other symbols
\else % if luatex or xetex
  \usepackage{unicode-math}
  \defaultfontfeatures{Scale=MatchLowercase}
  \defaultfontfeatures[\rmfamily]{Ligatures=TeX,Scale=1}
\fi
% Use upquote if available, for straight quotes in verbatim environments
\IfFileExists{upquote.sty}{\usepackage{upquote}}{}
\IfFileExists{microtype.sty}{% use microtype if available
  \usepackage[]{microtype}
  \UseMicrotypeSet[protrusion]{basicmath} % disable protrusion for tt fonts
}{}
\makeatletter
\@ifundefined{KOMAClassName}{% if non-KOMA class
  \IfFileExists{parskip.sty}{%
    \usepackage{parskip}
  }{% else
    \setlength{\parindent}{0pt}
    \setlength{\parskip}{6pt plus 2pt minus 1pt}}
}{% if KOMA class
  \KOMAoptions{parskip=half}}
\makeatother
\usepackage{xcolor}
\IfFileExists{xurl.sty}{\usepackage{xurl}}{} % add URL line breaks if available
\IfFileExists{bookmark.sty}{\usepackage{bookmark}}{\usepackage{hyperref}}
\hypersetup{
  hidelinks,
  pdfcreator={LaTeX via pandoc}}
\urlstyle{same} % disable monospaced font for URLs
\usepackage[margin=1in]{geometry}
\setlength{\emergencystretch}{3em} % prevent overfull lines
\providecommand{\tightlist}{%
  \setlength{\itemsep}{0pt}\setlength{\parskip}{0pt}}
\setcounter{secnumdepth}{-\maxdimen} % remove section numbering
\usepackage{esint}
\setmathfont{Asana Math}
\setlength{\emergencystretch}{3em}

\author{}
\date{}

\begin{document}

\hypertarget{removing-the-weyl-hypothesis-by-the-topology-of-ale-embeddings}{%
\section{Removing the Weyl hypothesis by the topology of ALE
embeddings}\label{removing-the-weyl-hypothesis-by-the-topology-of-ale-embeddings}}

24 September 2026. This note supplies the previously missing input in
the \href{20260923_compact_realization.md}{compact-regime realization
proof}. The stronger ALE nonembedding theorem is \textbf{an existing
result}, not a new classification discovered here: it is Zhang's theorem
with the noncompact extension stated explicitly by Bamler--Chen. Its
application changes the scope of the project's realization theorem
substantially.

In particular, the residual possibility recorded in
\href{20260923_topological_core_scope.md}{the preceding topology note},
namely a mixed-chirality rational-ball bubble in the universal cover,
can now be excluded. No classification of arbitrary Ricci-flat ALE
metrics is assumed.

\hypertarget{the-classical-input-that-closes-the-remaining-bubble-alternative}{%
\subsection{1. The classical input that closes the remaining bubble
alternative}\label{the-classical-input-that-closes-the-remaining-bubble-alternative}}

\textbf{ALE nonembedding lemma (Zhang; Bamler--Chen).} Let \(N\) be a
smooth four-manifold, possibly noncompact, with \[
 H_2(N;\mathbb Z)=0,\qquad H_1(N;\mathbb Z)\text{ torsion-free}.
 \tag{1}
\] No compact smooth domain \(B\subset N\) has interior diffeomorphic to
a nonflat complete Ricci-flat ALE four-manifold.

The precise noncompact formulation is
\href{https://arxiv.org/pdf/2305.03154}{Bamler--Chen, Lemma 4.35,
p.~41}, which cites
\href{https://msp.org/pjm/2019/303-1/pjm-v303-n1-p14-s.pdf}{Zhang,
Theorem 6.1 and Lemmas 6.2--6.5, pp.~378--382}. In the published version
of Bamler--Chen, the statement is still Lemma 4.35:
\href{https://doi.org/10.1002/cpa.70024}{DOI 10.1002/cpa.70024}. The
isolated \(H_1(M)\) in its printed statement denotes the ambient
\(H_1(N)\), as its following proof makes explicit. We use the lemma only
when this group is zero.

Here is the mechanism, to distinguish this input from the unavailable
classification of all Ricci-flat ALE spaces. In an ambient with
\(H_1=H_2=0\), write \(Y=\partial B=S^3/\Gamma\) and
\(C=N\setminus\operatorname{int}B\). Integral Mayer--Vietoris gives \[
 H_2(B;\mathbb Z)=0,\qquad
 H_1(Y;\mathbb Z)\simeq H_1(B;\mathbb Z)\oplus H_1(C;\mathbb Z).
 \tag{2}
\] The ALE fundamental group is finite, and the end group surjects onto
it. Thus \(H_1(B,Y;\mathbb Z)=0\). Duality and the universal coefficient
theorem identify the kernel of \(H_1(Y)\to H_1(B)\) with a group
isomorphic to \(H_1(B)\). The splitting in (2) respects this projection,
so \[
                  H_1(Y)\simeq G\oplus G.               \tag{3}
\] This is stronger than merely requiring that \(|H_1(Y)|\) be a square.

Zhang's spherical-space-form calculation reduces the nontrivial
possibilities in (3) to binary dihedral groups of order \(8n\),
\(n\ge1\), and the binary icosahedral group. His representation lemma
and eta calculation give, up to orientation, \[
 |\eta(Y)|=\frac{8n^2+1}{12n}\quad\text{or}\quad\frac{361}{180},
 \qquad\text{respectively}.                            \tag{4}
\] Both exceed \(2/3\). On the other hand, \(b_2(B)=0\), finite first
homology, connected boundary and duality give \(\chi(B)=1\),
\(\tau(B)=0\). Ricci-flat Chern--Gauss--Bonnet and APS signature, in
curvature-operator norms, imply \[
 \begin{split}
 E_++E_-&=8\pi^2\left(1-|\Gamma|^{-1}\right),\\
 E_+-E_-&=12\pi^2\eta(Y),\qquad
 E_\pm=\int|W^\pm|^2\ge0.
 \end{split}                                           \tag{5}
\] Consequently \[
 |\eta(Y)|\le\frac23\left(1-|\Gamma|^{-1}\right)<\frac23,
 \tag{6}
\] contradicting (4). The trivial end group gives zero total curvature
in (5), hence flatness. This explains why a mixed-chirality rational
ball does not evade the argument: nonnegativity of \textbf{both}
energies is already incompatible with its boundary data. The argument
does not set either Weyl half to zero.

For the quaternion test \(\Gamma=Q_8\), (3) is indeed possible: the
boundary has \(H_1=(\mathbb Z/2)^2\). It is (6), not this group
calculation alone, that excludes a Ricci-flat rational-ball filling:
\(3/4>7/12\). This separates the topological and geometric restrictions.

\hypertarget{vanishing-second-homotopy-suffices-torsion-in-the-first-is-allowed}{%
\subsection{2. Vanishing second homotopy suffices; torsion in the first
is
allowed}\label{vanishing-second-homotopy-suffices-torsion-in-the-first-is-allowed}}

\textbf{Theorem 1 (all-ALE-core exclusion).} Let \(M\) be a connected
smooth oriented four-manifold with \(\pi_2(M)=0\). No large compact core
of a nonflat complete smooth Ricci-flat ALE four-manifold embeds
smoothly in \(M\). No torsion-free assumption on \(\pi_1(M)\) is needed.

\textbf{Proof.} Let \(i:B\hookrightarrow M\) be such an embedding and
let \(p:\widetilde M\to M\) be the universal cover. Hurewicz gives \[
 H_1(\widetilde M;\mathbb Z)=H_2(\widetilde M;\mathbb Z)=0.
 \tag{7}
\] The fundamental group of a Ricci-flat ALE space is finite. A
component \(\widehat B\) of \(p^{-1}(i(B))\) is therefore a
\textbf{finite} cover of \(B\): its degree is the order of the finite
image of \(i_*:\pi_1(B)\to\pi_1(M)\). In particular, it is a compact
smooth domain in \(\widetilde M\). This step neither assumes that the
core fundamental group injects nor that its image is trivial.

The interior of \(\widehat B\) carries the lifted complete Ricci-flat
metric. A finite cover of an ALE space has finitely many lifted ALE ends
and Euclidean volume growth. It has only one end: if it had two, the
nonnegative-Ricci splitting theorem would split it as a line times a
compact manifold, contradicting Euclidean four-dimensional volume
growth. Thus the lifted space is itself nonflat Ricci-flat ALE. Apply
the classical lemma to its compact domain
\(\widehat B\subset\widetilde M\), using (7), to obtain a contradiction.
\(\square\)

For completeness, finiteness of the ALE fundamental group can also be
seen directly from the end: take the cover corresponding to the image of
\(\pi_1(Y)\to\pi_1(B)\). This subgroup is finite. If its index were
greater than one, one lifted end has degree one and is separated from
the rest by a compact cross-section. The complete lifted metric would
have at least two ends and Euclidean lower volume growth, again
contradicting splitting. Therefore the end group surjects onto the full
fundamental group. This reasoning uses Ricci-flatness and the ALE end,
not any holonomy classification. The same end-cover/splitting argument
appears in \href{https://arxiv.org/html/1102.2442}{Wright, Section 3 and
Proposition 1}; only that argument is used, not the paper's half-flat
classification.

The conclusion is strictly broader in topology than the previous
condition \(\pi_1\) torsion-free, \(\pi_2=0\). For example,
\(L(p,q)\times S^1\) has vanishing \(\pi_2\) and fundamental group
\((\mathbb Z/p)\times\mathbb Z\); homotopy four-spheres also have
vanishing \(\pi_2\). These illustrate the topological scope, not a claim
that every example has the near-Einstein inputs of the next theorem.

\hypertarget{bounded-geometry-realization-without-gauge-saturation-or-weyl-bounds}{%
\subsection{3. Bounded-geometry realization without gauge saturation or
Weyl
bounds}\label{bounded-geometry-realization-without-gauge-saturation-or-weyl-bounds}}

\textbf{Theorem 2 (harmonic compactness).} Fix \(K,D,v>0\) and
\(0<\alpha<1\). Closed connected smooth oriented four-manifolds with \[
 \pi_2(M)=0,\qquad
 \|\operatorname{Ric}^{\sharp}\|_{\infty,\mathrm{op}}\le K,
 \qquad\operatorname{diam}(g)\le D,\qquad V_g\ge v
 \tag{8}
\] have a uniformly positive \(C^{1,\alpha}\) harmonic radius. They are
precompact modulo diffeomorphisms in \(C^{1,\beta}\), \(\beta<1\), and
weak \(W^{2,p}\), \(p<\infty\), and have finitely many smooth types.

\textbf{Proof.} The analytic argument of
\href{20260923_compact_realization.md}{compact realization, Sections
3--4} is used up to, but not including, its one-half Weyl step.
Bishop--Gromov gives uniform noncollapse on fixed scales and a
total-volume upper bound. Jiang--Naber supplies the full
curvature-energy bound. If the minimum harmonic radius \(\rho_j\) tends
to zero, rescale by \(\rho_j^{-2}\). The normalized harmonic atlases
converge, Ricci tends uniformly to zero, and the harmonic-coordinate
metric equation upgrades convergence locally to strong \(W^{2,p}\) for
finite \(p\). The limit is smooth Ricci-flat. Recovery of harmonic
charts with strict margins excludes a flat limit.

Euclidean volume growth and finite curvature energy make this limit ALE
by Bando--Kasue--Nakajima. Only positive-order ALE asymptotics are
needed; no optimal decay-order assertion is used. Its smooth compact
core embeds in each sufficiently late \(M_j\), contradicting Theorem 1.
There is \textbf{no need to force \(W^+\) or \(W^-\) to vanish}. The
remaining finite-atlas compactness argument is unchanged. \(\square\)

\textbf{Theorem 3 (same-smooth-type realization).} Fix \(D,v>0\). There
is \(\varepsilon(D,v)>0\), uniform over smooth types and
\(\sigma\in\{-1,+1\}\), such that \[
 \begin{gathered}
 \pi_2(M)=0,\qquad \operatorname{diam}(g)\le D,\qquad V_g\ge v,\\
 \delta_\sigma(g):=
 \|\operatorname{Ric}^{\sharp}_g-\sigma I\|_{\infty,\mathrm{op}}
          <\varepsilon(D,v)
 \end{gathered}
 \quad\Longrightarrow\quad
 \exists h\text{ on }M:\ \operatorname{Ric}_h=\sigma h.
 \tag{9}
\]

\textbf{Proof.} Otherwise take counterexamples with the fixed data and
\(\delta_{\sigma_j}(g_j)\to0\). Pass to a subsequence with constant
sign. Theorem 2 applies with \(K=2\) to late terms and identifies them
with a fixed smooth manifold. In harmonic coordinates the limiting
metric satisfies \(\operatorname{Ric}=\sigma g\) weakly. Elliptic
regularity makes it smooth. Pull it back to the original smooth
manifolds, contradicting their defining absence of such an Einstein
metric. \(\square\)

In the positive case, \(\delta_+<1/2\) gives
\(\operatorname{Ric}\ge g/2\); Bonnet--Myers supplies
\(\operatorname{diam}\le\pi\sqrt6\). Thus for each \(v>0\), \(\pi_2=0\),
\(V_g\ge v\), and sufficiently small \(\delta_+\) already imply a
positive Einstein metric, with threshold depending only on \(v\). This
still retains the noncollapse input.

On the negative side an actual conformally solvable monopole class with
\((a_g^+)^2\ge A_0>0\) gives \[
 V_g\ge\frac{2\pi^2 A_0}{(1+\delta_-(g))^2}
       \ge\frac{8\pi^2 A_0}{9}\quad(\delta_-\le1/2).
 \tag{10}
\] Hence (9) yields a threshold \(\varepsilon(D,A_0)\) from this actual
class input and diameter alone. Unlike the preceding KE realization
theorem, \textbf{no smallness of \(\mathfrak D_2/V\) is required}. The
conclusion is negative Einstein, not necessarily Kähler--Einstein. That
already satisfies the existence output of the adopted Yau26 problem.

\hypertarget{what-is-resolved-and-what-is-not}{%
\subsection{4. What is resolved, and what is
not}\label{what-is-resolved-and-what-is-not}}

This closes the entire compact, noncollapsed bubble alternative on
\(\pi_2=0\) topology. A hypothetical negative-side Yau26 counterexample
sequence in that class, with uniform positive actual projection budget,
must have diameters tending to infinity: any bounded-diameter
subsequence would contradict (9)--(10). No near-saturation is necessary
for this escape conclusion.

On the positive side, if \(\pi_2(M_j)=0\), no \(M_j\) admits any
Einstein metric, and \(\delta_+(g_j)\to0\), then necessarily
\(V_{g_j}\to0\). Otherwise a subsequence has \(V_{g_j}\ge v>0\), and
Myers plus Theorem 3 produces a positive Einstein metric on every
sufficiently late member of that subsequence. Thus, within the
\(\pi_2=0\) topology, the two signs have different remaining escape
mechanisms: positive candidates must collapse in total volume; negative
candidates with a positive actual projection budget must escape in
diameter.

The original dimension-only problem still allows diameter escape,
collapse without a supplied volume budget, and topology with nonzero
\(\pi_2\). Those are not excluded here. The theorem does not decide
whether the fixed non-Kähler shear examples admit general Einstein
metrics: their known approximating sequence has diverging diameter.

The old near-saturation theorem retains its independent stronger
geometric output: its actual forms recover a global parallel form and
hence KE. For \textbf{general Einstein existence} in the compact regime,
however, saturation and the one-sided Weyl assumptions are now
unnecessary.

The mathematical increment in this note is the explicit composition of
the full ALE nonembedding input, the finite-cover lift, and the
preserved almost-Einstein compactness argument. Zhang's boundary/eta
obstruction and Bamler--Chen's noncompact formulation are credited as
existing results. No assertion of world-first priority for the resulting
compact-regime corollary, or classification of all Ricci-flat ALE
spaces, is made.

For the historical comparison,
\href{https://www.math.stonybrook.edu/~anderson/einsurv.pdf}{Anderson's
survey, Corollary 5.8} already gives exclusion of the orbifold boundary
of the exact Einstein moduli space under vanishing \(H_2(M;\mathbf F)\)
for every field. Bamler--Chen explicitly credit Anderson's earlier work
for the nonembedding idea. The composition here uses the later explicit
noncompact lemma, the finite-cover lift under \(\pi_2=0\), and
bounded-Ricci harmonic compactness to handle varying smooth types and
pointwise approximate Einstein inputs. This distinguishes the statements
being used; it does not establish that this composed corollary was
absent from the literature.

\end{document}
